\documentclass[11pt,a4paper]{ctexart}
\usepackage[margin=2.45cm,marginparwidth=1.15cm]{geometry}
\usepackage{amsmath,amssymb,amsthm,mathtools,bm}
\usepackage{graphicx}
\usepackage{xcolor}
\usepackage{fancyhdr}
\usepackage[colorlinks=true,linkcolor=blue!45!black,urlcolor=blue!45!black]{hyperref}
\usepackage{enumitem}
\setlength{\parindent}{2em}
\linespread{1.28}
\pagestyle{fancy}
\fancyhf{}
\fancyfoot[C]{\thepage}
\renewcommand{\headrulewidth}{0pt}
\newcommand{\SourcePage}[1]{\marginpar{\raggedright\footnotesize\color{gray}原文 p.#1}}
\newcommand{\TN}[1]{\footnote{[译注]\,#1}}
\newcommand{\term}[1]{\textsf{#1}}
\title{\bfseries 弦论中的 $K$-理论应用概览}
\author{Edward Witten\\[2pt]
\small School of Natural Sciences, Institute for Advanced Study\\[-2pt]
\small 原载 arXiv:hep-th/0007175（2000 年 8 月）}
\date{中文译注本 \; v2.0 重制版（2026-11-14）}
\begin{document}
\maketitle
\thispagestyle{empty}

\begin{center}
\begin{minipage}{0.86\textwidth}
\small
\textbf{重制说明.} 本册按《数学文献翻译纪律》v2.0 重制：以原文为底本逐段核对，公式、编号、文献与页锚点均回查原文；旧版（Markdown$\to$DOCX$\to$PDF）已隔离。
译注以脚注标出，凡属译者补充而原文所无者，均注明依据或标「译者按（无外部依据）」。
\end{minipage}
\end{center}

\vspace{0.6em}
\(K\)-理论为分类 Ramond-Ramond (RR) 荷与 RR 场提供了一个框架。流形的 \(K\)-理论可以自然地推广为非交换代数的 \(K\)-理论，例如非交换 Yang-Mills 理论或开弦场论中所考虑的那些代数。在许多具体问题中，如果从无穷多个 \(D\)-膜出发、再经由快子凝聚约化为有限多个，那么 \(K\)-理论的分析会进行得最为自然。这提示我们，应当对 \(N=\infty\) 的情形重新考虑弦场论。

\section{\texorpdfstring{引言与 \(K(X)\) 的定义}{引言与 K(X) 的定义}}

\SourcePage{2}包裹在时空子流形 \(S\) 上的 \(D\)-膜可以携带非零的 Ramond-Ramond (RR) 荷。RR 场是 \(p\)-形式，表面看来，守恒荷似乎应当由 RR 形式（或由循环 \(S\) 本身）的上同调类来度量。然而，\(D\)-膜携带规范场；而规范场在上同调理论中并不自然。它们在“\(K\)-理论”中才是自然的。\(K\)-理论已被用于回答关于 RR 荷与 RR 场的一些问题；本文的目的是给出这方面的概述，并附上一些猜测（我的证据非常有限）：关于在大 \(N\) 极限下应当如何重新思考开弦场论。这些内容构成第 1-3 节。一些数学细节已推迟到第 4 节。（本文基于我在 2000 年 7 月 10 日于密歇根大学 Strings 2000 的演讲，以及在 CIT-USC 理论物理中心更早的若干演讲。我感谢两个听众的提问与评论。）

若 \(X\) 是时空，\(\mathcal{A}(X)\) 是 \(X\) 上连续复值函数的交换结合代数，则 \(X\) 的 \(K\)-理论可以用 \(\mathcal{A}(X)\) 的表示来定义。环的一个表示通常称为模。下面是 \(\mathcal{A}(X)\)-模的一些例子。

最显然的 \(\mathcal{A}(X)\)-模例子就是 \(\mathcal{A}(X)\) 本身。对 \(f\in  \mathcal{A}(X)\)（视为环）与 \(g\in \mathcal{A}(X)\)（视为模），我们定义

\begin{equation}\tag{1.1}\label{eq:1.1}
f(g)=fg,
\end{equation}

其中右端的乘法在 \(\mathcal{A}(X)\) 中进行。这显然满足模的定义条件，即 \((f_1f_2)(g)=f_1(f_2(g))\)。

更一般地，考虑包裹在 \(X\) 的子流形 \(S\) 上的一个 \(Dp\)-膜（或某个 \(N>0\) 的 \(N\) 个 \(Dp\)-膜的集合），并带有 \(D\)-膜上任意的 Chan-Paton 规范丛 \(W\)。令 \(M(S)\) 为 \(W\) 的截面空间，即与丛 \(W\) 耦合的带电标量的单粒子态空间。于是 \(M(S)\) 是一个 \(\mathcal{A}(X)\)-模；对 \(f\in \mathcal{A}(X)\)、\(g\in M(S)\)，我们同样置 \(f(g)=fg\)。在右端，乘法定义为把 \(f\)（它是 \(X\) 上的函数）限制到 \(S\) 上，然后将 \(f\) 与 \(g\) 相乘。

于是，比如说在 IIB 型超弦理论中，一组 \(D9\)-膜定义了 \(\mathcal{A}(X)\) 的一个表示或模 \(E\)。一组 \(\overline{D9}\)-膜定义了另一个模 \(F\)。因此任何 \(D9\) 与 \(\overline{D9}\)-膜的构型都确定一个对 \((E,F)\)。

\SourcePage{3}为了分类 \(D\)-膜荷，我们希望按物理过程来分类这些对 \((E,F)\)。一个重要的过程（Sen 1998）是膜-反膜对的产生与湮灭——即产生或湮灭一组各自带有相同规范丛 \(G\) 的 \(D9\) 与 \(\overline{D9}\)。这相当于

\begin{equation}\tag{1.2}\label{eq:1.2}
(E,F)\leftrightarrow(E\oplus G,F\oplus G).
\end{equation}

这些等价类构成一个群，称为 \(K(X)\)（若想更明确地体现用 \(\mathcal{A}(X)\)-模作出的解释，则记作 \(K(\mathcal{A}(X))\)）。该群的加法就是

\begin{equation}\tag{1.3}\label{eq:1.3}
(E,F)+(E',F')=(E\oplus E',F\oplus F').
\end{equation}

\((E,F)\) 的逆是 \((F,E)\)；注意 \((E,F)\oplus (F,E)=(E\oplus F,E\oplus F)\)，而利用等价关系（上式），它等价于零。IIB 型的 \(D\)-膜携带在 \(K(X)\) 中取值的守恒荷。

在上述 \(K(X)\) 的定义中（Witten 1998），我们只用了九维膜，尽管如前面所解释的，对任意 \(p\) 都可以用 \(Dp\)-膜（或反膜）构造一个 \(\mathcal{A}(X)\) 模。事实上，我们可以只用九维膜来分类 \(D\)-膜荷，然后通过带有适当快子凝聚的对 \((E,F)\) 构造出 \(p<9\) 的 \(Dp\)-膜。这一构造（源于 Atiyah、Bott 和 Shapiro）已在（Witten 1998）第 4 节中、于 \(D\)-膜背景下作了回顾。关于用模来定义 \(K(X)\)、并阐明九维膜所起作用的更系统解释，可在第 4 节找到。

\subsection{\texorpdfstring{\(K\)-理论描述的优点}{K-理论描述的优点}}

知道 \(D\)-膜荷由 \(K\)-理论分类，我们能得到什么？

首先，这是正确的答案。凡是仔细考察 RR 荷（或 RR 场）拓扑性质的地方，都会看到反映 \(K\)-理论结构的效应。例如，存在稳定的 \(D\)-膜态（如 I 型理论中非超对称的 \(D0\)-膜），如果 \(D\)-膜荷由上同调而非 \(K\)-理论分类，这些态就不会存在。反过来，也可能存在这样一种 \(D\)-膜态：若用上同调来度量 \(D\)-膜荷它是稳定的，但实际上并不稳定（通过一个涉及在中间态中成核产生 \(D9\)-\(\overline{D9}\) 对的过程）。这种情况在（Diaconescu、Moore 和 Witten 2000）II 型超弦理论中会发生：在其中，包裹在时空里一个同调非平凡循环上的 \(D\)-膜在某些情形下实际上是不稳定的。最后，\(D\)-膜的 \(K\)-理论解释对于理解某种整体世界体反常\TN{原文作 global worldvolume anomaly；旧译「整体世界体异常」，此处按物理学惯例改「反常」。}是必需的（Freed 和 Witten 1999）。

但我认为，知道 \(D\)-膜的 \(K\)-理论解释还有一个更深刻的理由：它可能自然地适用于弦论式的推广。\SourcePage{4}事实上（尽管一些数学细节已推迟到第 4 节），我们是用时空中函数代数 \(\mathcal{A}(X)\) 的表示来定义 \(K(X)\) 的。我们同样可以用 \(\mathcal{A}\)-模的对 \((E,F)\)，对任何非交换代数 \(\mathcal{A}\) 定义 \(K(\mathcal{A})\)。相比之下，如果把时空上的函数代数换成非交换环，我们就不会有同样有用而方便的“上同调”概念。

例如，开启一个 Neveu-Schwarz \(B\)-场，我们可以使 \(\mathcal{A}(X)\) 成为非交换的；相应的 \(K(\mathcal{A})\) 被 Connes、Douglas 和 Schwarz 在将非交换 Yang-Mills 理论应用于弦论的最初那篇论文中使用过（Connes、Douglas 和 Schwarz 1998）。这是一个有趣的例子，尽管它只涉及弦的零模。人们非常希望有一个完全的弦论版本，其中所用的非交换代数由弦的全部模式而不只是零模构造。

那么，使用全部模式的正确非交换代数是什么？我们当然不知道。一个具体的候选者是开弦场论的 \(*\)-代数，它用把弦粘合在一起的方式来定义。如果我把这个代数记作 \(\mathcal{A}_{st}\)，那么 \(D\)-膜荷由 \(K(\mathcal{A}_{st})\) 自然地标记，看来是合理的。（对于体积远大于弦标度的流形，我猜测 \(K(\mathcal{A}_{st})\) 与拓扑 \(K\)-理论的普通 \(K(X)\) 相同。）我将在第 3.3 节回到 \(K(\mathcal{A}_{st})\)。首先，我想完成对 \(K\)-理论在弦物理中已知应用的综述。

\section{\texorpdfstring{\(K\)-理论与 RR 场}{K-理论与 RR 场}}

\(K\)-理论不仅与理解 RR 荷有关，也与理解 RR 场有关 [6--8,3]。

粗略地说，RR 的 \(p\)-形式场 \(G_p\) 服从 Dirac 量子化律：对时空中任一 \(p\)-循环 \(U\)，

\begin{equation}\tag{2.1}\label{eq:2.1}
\int_U\frac{G_p}{2\pi}=\,{\mathrm integer}.
\end{equation}

如果这就是正确的条件，那么 RR 场就会由上同调分类。

但这不是正确的答案，因为 RR 周期上实际的量子化条件比\eqref{eq:2.1}微妙得多。对\eqref{eq:2.1}存在各种各样的修正，它们涉及时空曲率与膜上的规范场，也涉及自对偶性与整体异常 [9--13,7,8,3]。

对 IIB 型超弦，答案原来是：RR 场由 \(K^1(X)\) 分类。就我们的目的而言，\(K^1(X)\) 可以定义为从 \(X\) 到 \(U(N)\) 的连续映射群的分量群，其中 \(N\) 任意足够大。

\SourcePage{5}这一陈述意味着，\(X\) 上 RR 场的拓扑类由某个大 \(N\) 的映射 \(U:X\to U(N)\) 分类。\(G_p\) 与 \(U\) 的关系大致是 \(G_p\sim \mathrm{Tr}\,(U^{-1}dU)^p\)；这里我忽略了时空曲率引起的修正，以及与 RR 场自对偶性相关的微妙之处。

\(U\) 的物理意义并不清楚。对 IIA 型，类似的说法是：RR 场由一个带联络 \(A\) 与曲率 \(F_A\) 的 \(U(N)\) 规范丛（某个大 \(N\)）分类，关系为 \(G_p\sim \mathrm{Tr}\,F_A^{p/2}\)。对 \(M\)-理论，类似物涉及带联络的 \(E_8\) 规范丛 [13,14,3]。同样，这些 \(U(N)\) 或 \(E_8\) 规范场的物理意义并不清楚。

用 \(K^1\) 来分类 IIB 型 RR 场的价值在于，它给出了一种简洁的方式来概括 RR 场所服从的、本来相当复杂的量子化条件。此外，这一框架有助于描述进入 RR 配分函数的微妙相位因子。事后看来，一旦知道 RR 荷由 \(K\)-理论分类，人们就应当猜测 RR 场也有类似的分类。毕竟，RR 荷产生 RR 场！所以用于分类 RR 荷的数学，必然与用于分类 RR 场的数学相似。

正如 \(K(X)\) 那样，\(K^1(X)\) 对任何非交换代数 \(\mathcal{A}\) 都有类似物。给定 \(\mathcal{A}\)，我们以 \(\mathcal{A}_N\) 记其矩阵元属于 \(\mathcal{A}\) 的可逆 \(N\times N\) 矩阵群。那么 \(K^1(\mathcal{A})\) 就是 \(\mathcal{A}_N\) 对大的 \(N\) 的分量群。

例如，当 \(\mathcal{A}=\mathcal{A}(X)\) 为 \(X\) 上复值连续函数环时，\(\mathcal{A}_N\) 是从 \(X\) 到 \(GL(N,{\mathbf C})\) 的映射群。它可以收缩到从 \(X\) 到 \(U(N)\) 的映射群，因此对大的 \(N\)，\(\mathcal{A}_N\) 的分量群与我们最初定义的 \(K^1(X)\) 相同。

\(K^1(X)\) 对非交换环的推广存在，这意味着用 \(K^1(X)\) 对长距离极限下 IIB 型 RR 场所作的描述，可能成为弦论式推广的有用起点。

\section{\texorpdfstring{\(N\to\infty\)}{N 趋于无穷}}

上一节中，\(N\) 是一个足够大但取有限的整数。我们接下来的任务，是描述一些依赖于把 \(N\) 取为无穷的事情。

在这样做之前，让我们回顾一下 \(N\to\infty\) 极限在物理学中的角色。它在规范理论与弦之间猜想性的联系中很重要；在用于给出弦论可解例子的旧矩阵模型中很重要；在 \(M\)-理论的矩阵模型中很重要；\SourcePage{6}也在渐近 AdS 时空中的引力与边界上的共形场论之间的对应中很重要。

我这里的主题，是想提议：我们应当以某种方式研究 \(N=\infty\) 的 \(D\)-膜，并借助快子凝聚湮灭掉绝大部分膜，从而约化到某种更易于把握的东西。为说明这样做的动机，我将考虑两个具体问题，它们看来都要求把 \(D\)-膜的数目取为无穷。一个问题是 Type IIA 超弦与 \(K\)-理论的关系；另一个是把拓扑非平凡的 NS 三形式场 \(H\) 纳入 RR 荷与 RR 场的 \(K\)-理论分类。

在作完报告之后，我才知道有一个标准问题，早已被用无穷多个 \(D9\) 与 \(\overline{D9}\)-膜以及快子凝聚到某种可处理的东西的方式解释过了（见 Hori 1999，第 2.2、4、5 节）。这是一个涉及 \(D5\)-\(D9\) 系统的 \(D5\)-膜探针的问题。

\subsection{Type IIA}

对于 Type IIB 超弦，我们用 \(K(X)\) 分类 RR 荷，用 \(K^1(X)\) 分类 RR 场。

\(T\)-对偶的说法则是：对于 Type IIA，应当用 \(K^1(X)\) 分类 RR 荷，而用 \(K(X)\) 分类 RR 场。（由 Bott 周期性，\(K^{i+2}(X)=K^i(X)\)，因此 \(X\) 的 \(K\)-群只有 \(K^0(X)\)——我们一直简称为 \(K(X)\)——和 \(K^1(X)\)。）

想要一般地解释 Type IIA 中为何由 \(K^1(X)\) 分类 RR 荷，最具体也最自然的尝试是 Horava 的（Horava 1999）。这里的出发点是考虑 Type IIA 的 \(N\) 个不稳定 \(D9\)-膜组成的系统。这些膜上支持一个 \(U(N)\) 规范场，以及一个处于 \(U(N)\) 伴随表示中的快子场 \(T\)。存在一个对称性 \(T\to -T\)。

快子场的有效势被认为具有如下一般形式 \begin{equation}\tag{3.1}\label{eq:3.1}
V(T)=\frac{1}{g_{st}}\mathrm{Tr}\,F(T),
\end{equation} 其中函数 \(F(T)\) 非负，并且在把 \(T\) 作适当标度之后，恰好当 \(T=\pm 1\) 时为零。因此 \(V(T)\) 取极小，当且仅当 \(T\) 的每个本征值都是 \(\pm 1\)。

文献（Horava 1999）中论证了：在平坦的 \(\mathbb{R}^{10}\) 中，可以把超对称的 \(Dp\)-膜（\(p\) 为偶数）构造为 \(T\) 的孤子。例如，要造一个 \(D6\)-膜，我们取 \(N=2\)。令 \(\vec x\) 为垂直于该 \(D6\)-膜的三个坐标，并令 \begin{equation}\tag{3.2}\label{eq:3.2}
T=\frac{\vec\sigma\cdot\vec x}{|x|}f(|x|)
\end{equation} \SourcePage{7}其中当 \(r\) 很小时 \(f(r)\sim r\)，而当 \(r\to\infty\) 时 \(f(r)\to\infty\)。于是在 \(|x|\to\infty\) 时，\(T\) 的本征值处处为 \(\pm 1\)。在 \(x=0\) 附近存在一个拓扑性的“结”\TN{原文 topological knot，指快子场本征值在无穷远处的缠绕，非纽结理论中的「结」。}，我们把它解释为该 \(D6\)-膜。

在平坦的 \({\mathbf R}^{10}\) 中，人们同样可以为其它的偶数 \(p\) 造出 \(Dp\)-膜。但在一般的时空上，除非取 \(N=\infty\)，否则这对任意的 \(Dp\)-膜都行不通。如果 \(X\)，或至少它的空间部分，是紧致的，这个问题就最为明显。快子场 \(T\) 取伴随表示的值，因而把 \(X\) 映射到 \(U(N)\) 的李代数中；由于李代数可缩，\(T\) 不携带任何拓扑。所以从 \(X\) 到李代数的映射并不代表 \(K^1(X)\) 的一个元素；事实上，它根本不携带拓扑信息。

要定义 \(K^1(X)\) 的一个元素，我们需要的是群而不是李代数；一个映射 \(U:X\to U(N)\) 就能做到这一点，正如我前面已经说过的。

令人惊讶的是，正如 Atiyah 和 Singer 很久以前所证明的，只要取 \(N=\infty\)，我们就能从李代数重新得到正确的拓扑！我们必须把 \(U(\infty)\) 解释为一个可数无穷维 Hilbert 空间 \({\mathcal H}\) 上的幺正群 \({\mathcal U}\)。（这样的 Hilbert 空间也称为可分 Hilbert 空间。）我们把厄米 \(N\times N\) 矩阵所成空间的 \(N=\infty\) 类比，解释为 \({\mathcal H}\) 上有界自伴算子 \(T\) 所成的空间，其谱如下：有无穷多个正本征值和无穷多个负本征值，且零不是谱的聚点。（最后一个条件使 \(T\) 成为“Fredholm 算子”。\TN{原文 Fredholm operator；此处指有界自伴、谱离散且零非聚点的算子。}）从物理上说，应当要求 \(T\) 满足这些条件，因为只有这样才能使能量与 \(D8\)-膜荷保持有限。事实上，为使能量有限，\(T\) 的几乎所有本征值都非常接近 \(\pm 1\)。无论如何，在把这些条件加于 \(T\) 之后，\(T\) 所成的空间与 \(N\) 很大时 \(U(N)\) 的空间具有相同的拓扑。

于是我们可以用 \(D9\)-膜系统上的快子凝聚来描述 Type IIA 的 RR 荷。但我们必须从无穷多个 \(D9\)-膜出发，它们随后经过快子凝聚，落到一个有限能量的构型上。

\textbf{有限 \(N\) 的障碍}

现在让我们更具体地说明：用这种方式在有限 \(N\) 下构造 \(D\)-膜的障碍是什么，以及它为何在 \(N=\infty\) 时消失。

让我们回到用 2 个 \(D9\)-膜构造一个 \(D6\)-膜的例子。我们把横向方向取为一份 \(\mathbb{R}^3\)，并把快子场取为 \begin{equation}\tag{3.3}\label{eq:3.3}
T=\frac{\vec\sigma\cdot\vec x}{|x|}f(|x|)
\end{equation} \SourcePage{8}如果我们试图把横向方向紧致化成 \({\mathbf S}^3\)，就会遇到麻烦，因为 \(T\) 在无穷远处不是常数。\(T\) 的共轭类在无穷远处是常数——\(T\) 的本征值处处为 \(1\) 和 \(-1\)——但 \(T\) 本身不是常数。

此外，\(T\) 在无穷远处不同伦于一个常数。如果 \(T\) 在无穷远附近同伦于常数，我们就可以把它形变为常数，然后延拓到整个 \({\mathbf S}^3\) 上。但它并不同伦于常数。

使 \(T\) 在无穷远处成为常数的基本障碍是“磁荷”。设 \({\mathbf S}^2\) 是 \(\mathbb{R}^3\) 中无穷远处的一个球面。在 \({\mathbf S}^2\) 上，我们可以定义一个线丛 \({\mathcal L}_+\)，其纤维是 \(T\) 的 \(+1\) 本征空间；再定义一个线丛 \({\mathcal L}_-\)，其纤维是 \(T\) 的 \(-1\) 本征空间。这两个线丛 \({\mathcal L}_+\) 与 \({\mathcal L}_-\) 在拓扑上是非平凡的——它们的第一 Chern 类分别为 1 和 \(-1\)。只要我们试图在保持 \(T\) 的本征值为 1 和 \(-1\) 的前提下形变 \(T\)，线丛 \({\mathcal L}_\pm\) 就是良定义的，它们的第一 Chern 类也不变。因此 \({\mathcal L}_+\)（或 \({\mathcal L}_-\)）的非平凡性阻止我们把 \(T\) 同伦地形变为常数。

让我们再加上一些额外的“旁观” \(D9\)-膜，看看会有什么变化。设有 \(M=2k\) 个额外的膜，于是膜的总数为 \(N=2+M=2+2k\)。令快子场为 \(T'=T\oplus U\)，其中 \(T\) 如上，而 \(U\) 是作用在这 \(2k\) 个额外膜上的 \(k\) 份矩阵 \begin{equation}\tag{3.4}\label{eq:3.4}
\left(\begin{matrix}1 & 0 \\ 0 & -1\\\end{matrix}\right)
\end{equation} 之和。

于是 \(T'\) 在无穷远附近有 \(k+1\) 个本征值 \(+1\) 和 \(k+1\) 个本征值 \(-1\)。\(T'\) 的 \(+1\) 本征空间是一个第一 Chern 类为 \(1\) 的向量丛 \(V_+\)（因为它是在 \({\mathcal L}_+\) 上加一个平凡丛构造出来的），而 \(T'\) 的 \(-1\) 本征空间同样是一个第一 Chern 类为 \(-1\) 的向量丛 \(V_-\)。特别地，\(V_+\) 与 \(V_-\) 都是非平凡的，所以加上这些旁观膜并没有给我们带来任何好处：\(T'\) 不同伦于常数，也不能延拓到无穷远处之外。\(V_+\) 的非平凡性由 \(\pi_1(U(k+1))=\mathbb{Z}\) 控制，而后者与第一 Chern 类的存在相关联。

反之，如果我们取 \(k=\infty\) 会怎样？更确切地说，我们取旁观 \(D9\)-膜的数目为可数无穷，并假定 \(T'=T\oplus U\)，其中 \(U\) 是\eqref{eq:3.4}中矩阵的可数无穷多份拷贝的直和。我们仍然可以定义丛 \(V_+\) 与 \(V_-\)；它们的纤维是可分 Hilbert 空间（即可数无穷维的 Hilbert 空间）。此时 \(U(k+1)\) 被替换为 \({\mathcal U}\)，即一个可分 Hilbert 空间的幺正群。于是我们遇到一个基本事实（Kuiper 定理）：\({\mathcal U}\) 是可缩的；它的所有\SourcePage{9}同伦群都为零。因此任何以可分 Hilbert 空间为纤维的丛都是平凡的。特别地，\(V_+\) 与 \(V_-\) 都是平凡的，于是 \(T'\) 同伦于常数，并且可以延拓到无穷远处之外。

所以，如果不稳定 \(D9\)-膜的总数为 \(N=\infty\)，我们就能造出一个定域在 \({\mathbf S}^3\) 中某点上的 \(D6\)-膜。更一般地，鉴于 Atiyah 与 Singer 的结果，我们可以从 \(N=\infty\) 出发，通过快子凝聚构造出 \(K^1(X)\) 中的任意一个类。

就把这一结果应用于物理而言，有几个问题是我们应当担心的。一个问题是：从无穷多个膜出发、并依靠快子凝聚把我们带到有限能量的东西，这在物理上是否合理。我们只能把这个问题留给未来（但可参见 Hori 1999 中的一个例子）。

一个颇为不同的问题是：我们通过取 \(N=\infty\) 所得到的答案，对物理来说是否是正确答案。在定域于 \({\mathbf S}^3\) 上某点的 \(D6\)-膜场中，RR 二形式场 \(G_2\)（\(D6\)-膜是它的磁源）的方程没有解，因为“磁通无处可去”。

看来情况是这样的：\(N=\infty\) 对应于经典开弦理论中的正确答案，在那里有效作用量来自拓扑为圆盘的世界面。RR 场作为相对阶为 \(g_{st}\)（闭弦耦合常数）的修正进入，来自拓扑为圆柱的世界面，在经典近似中应当被忽略。

用膜的创生与湮灭来分类 \(D\)-膜，在 \(g_{st}=0\) 时成立，并且（正如我们在第 1 节中所回顾的）对 Type IIB 导致用 \(K(X)\) 分类 \(D\)-膜荷。要通过不稳定 \(D9\)-膜与快子凝聚得到 Type IIA 的类似答案——即 \(K^1(X)\)——我们需要从 \(N=\infty\) 出发。

直观地说，在没有快子凝聚时，\(N=\infty\) 应当对应于 \(g_{st}=0\)，因为开弦的有效展开参数是 \(g_{st}N\)。如果 \(N\) 是无穷，那么在快子凝聚之前，\(g_{st}\) 必须为零，否则量子修正会发散。如果我们希望 \(g_{st}\) 非零，就需要靠快子凝聚约化到一个有效的有限 \(N\) 值。

\subsection{\texorpdfstring{打开 \(H\) 场}{打开 H 场}}

一个有些类似的问题是：当 Neveu-Schwarz 三形式场 \(H\) 在拓扑上非平凡时，如何考虑 \(D\)-膜。我们将在 Type IIB 中展开这一讨论（对 Type IIA，我们须把下面的内容与上面关于 \(H\) 场不存在时所说的内容结合起来）。

正如 \(H=0\) \SourcePage{10}时一样，我们希望用满足通常那类等价关系的对 \((E,F)\)（其中 \(E\) 是一个 \(D9\) 态，\(F\) 是一个 \(\overline{D9}\) 态）来分类 \(D\)-膜态。但在有 \(H\) 场存在时，要有一个 \(D9\) 态是有问题的。

事实上，当 \(H\) 在拓扑上非平凡时，不可能有单个 \(D9\)-膜。在 \(D9\)-膜上存在一个场强为 \(F\) 的 \(U(1)\) 规范场。关系式 \(dF=H\) 表明，在 de Rham 上同调的意义上，若存在单个 \(D9\)-膜，则 \(H\) 必须是拓扑平凡的。这个结论其实是精确成立的，而不只是在 de Rham 上同调中成立。

有一种特殊情形，对这一困难有一个相对初等的补救办法（见 Witten 2000，第 5.3 节，以及 Kapustin 1999）。如果 \(H\) 是挠的，即存在整数 \(M>0\) 使 \(MH\) 拓扑平凡，那么就有可能有一组 \(M\) 个 \(D9\)-膜，其“规范丛”的结构群实际上是 \(U(M)/\mathbb{Z}_M\) 而不是 \(U(M)\)。（把 \(U(M)/\mathbb{Z}_M\) 丛提升为 \(U(M)\) 丛的障碍由 \(H\) 决定。）我们把这样的规范丛称为扭曲丛。更一般地，对任何正整数 \(m\)，我们可以有 \(N=mM\) 个 \(D9\)-膜，其丛的结构群为 \(U(mM)/\mathbb{Z}_M\)。在这种情况下，正如人们所猜测的，\(D\)-膜荷由满足通常等价关系的扭曲丛（或 \(D9\) 与 \(\overline{D9}\) 态）的对 \((E,F)\) 来分类。其等价类构成一个群 \(K_H(X)\)。

如果想把 \(K_H(X)\) 解释为某个代数的表示的 \(K\)-理论，就必须挑定一个特定的扭曲丛 \(W\)，并考虑边界条件由 \(W\) 决定的 \(D\)-膜态。\(W\)-\(W\) 开弦在伴随表示中变换，因此开弦零模部分的规范参数是 \(W\otimes \overline W\) 的截面。注意，尽管 \(W\) 是扭曲丛（结构群为 \(U(M)/\mathbb{Z}_M\) 而不是 \(U(M)\)），\(W\otimes \overline W\) 却是通常的丛，因为中心在伴随表示中作用平凡。

\(W\otimes \overline W\) 的截面构成一个代数，其定义如下：若 \(s^i{}_j\) 与 \(t^k{}_l\) 是 \(W\otimes \overline W\) 的截面——其中上下指标分别取 \(W\) 与 \(\overline W\) 中的值——则它们的乘积为 \((st)^i{}_l=\sum_ks^i{}_kt^k{}_l\)。这就是丛 \(W\) 的全部自同态或线性变换所成的代数 \({\mathcal A}_W(X)\)。对于 \(W\)-\(W\) 开弦，开弦场论的代数在只考虑弦的零模时就约化为 \({\mathcal A}_W\)。在 \(X\) 远大于弦标度的极限下，这是低能下的一个合理近似。

如果 \(H\) 为零，且 \(W\) \SourcePage{11}是平凡的秩一复丛，那么 \({\mathcal A}_W(X)\) 就是我们的老朋友 \({\mathcal A}(X)\)。如果 \(H\) 为零，且 \(W\) 是平凡的秩 \(N\) 复丛，那么引入 \(W\) 只意味着处处都有 \(N\times N\) 的 Chan-Paton 矩阵。所以在此情形下 \({\mathcal A}_W(X)={\mathcal A}(X)\otimes M_N\)，其中 \(M_N\) 是 \(N\times N\) 复值矩阵的代数。一般地，无论 \(H\) 如何，\(W\) 局部总是平凡的，因此局部上 \({\mathcal A}_W(X)\) 同构于 \({\mathcal A}(X)\otimes M_N\)。

一个扭曲丛等价于一个 \({\mathcal A}_W\)-模，而扭曲丛的对 \((E,F)\)（模去通常的等价关系）所成的群 \(K_H(X)\) 与 \(K({\mathcal A}_W)\)——\({\mathcal A}_W\)-模的 \(K\)-群——重合。这一点将在第 4 节中解释。这个论断立刻带来一个疑难：用扭曲丛的对 \((E,F)\) 定义的 \(K_H(X)\) 显然与 \(W\) 无关，而 \(K({\mathcal A}_W)\) 看起来却依赖于 \(W\)。的确，任给两个不同的扭曲丛 \(W\) 与 \(W'\)，相应的代数 \({\mathcal A}_W\) 与 \({\mathcal A}_{W'}\) 是不同的，但彼此“Morita 等价”。这一概念将在第 4 节解释，那里我们还会证明 Morita 等价蕴涵 \(K({\mathcal A}_W)=K({\mathcal A}_{W'})\)。

到目前为止，我们只考虑了 \(H\) 为挠的情形。一个在 AdS/CFT 对应中很重要的典型例子是时空 \(X={\mathrm AdS}_5\times {\mathrm RP}^5\)，其中 \({\mathrm RP}^5\) 上的一个挠 \(H\) 场被用来在边界共形场论中描述 \(Sp(n)\) 而非 \(SO(2n)\) 规范理论。

然而，在大多数物理应用中，\(H\) 不是挠的。在这种情况下，我们必须以某种方式取上面所述内容的大 \(M\) 极限。做到这一点的正确方式已由 Bouwknegt 与 Mathai（Bouwknegt 与 Mathai 2000）以及 Atiyah 与 Segal（Atiyah 与 Segal，未发表）给出。\(U(M)/\mathbb{Z}_M\) 的合适大 \(M\) 极限是 \(PU({\mathcal H})=U({\mathcal H})/U(1)\)。换言之，对 \(M=\infty\)，我们把 \(U(M)\) 替换为可分 Hilbert 空间 \({\mathcal H}\) 的幺正群 \(U({\mathcal H})\)，并把 \(\mathbb{Z}_M\) 替换为 \(U(1)\)。这特别意味着，当 \(H\) 不是挠的时候，不可能有有限个 \(D9\)- 或 \(\overline{D9}\)-膜，但可以有无限多个，并带有合适的无穷秩扭曲规范丛 \(E\) 或 \(F\)。此时 \(D\)-膜荷由对 \((E,F)\) 模去通常等价关系所成的群 \(K_H\) 来分类。详细的解释可在（Bouwknegt 与 Mathai 2000）中找到。Kuiper 定理——\({\mathcal U}=U({\mathcal H})\) 的可缩性——起着重要作用，正如它在第 3.1 节中那样。

这一构造在 \(M=\infty\) 极限下具有一个美妙的性质（见于（Bouwknegt 与 Mathai 2000））：其 \(K\)-群为 \(K_H\) 的那个非交换代数是\textbf{唯一}的，与扭曲丛 \(W\) 或 \(W'\) 的任何任意选取无关。这一点真正依赖于 \(D9\) 与 \(\overline{D9}\) 膜的数目为无穷。

\SourcePage{12}\subsection{弦论的推广？}

冒着稍稍冒险的风险，我现在想尝试给出其中一些内容的某种弦论式推广。

我们从一个闭弦背景出发，试图构造一个开弦代数。为此，我们必须挑定一个开弦边界条件——称之为 \(\alpha\)。于是在开弦场论中，\(\alpha\)-\(\alpha\) 开弦构成一个代数 \({\mathcal A}_\alpha\)。

这个代数并不唯一；我们也可以挑另一个边界条件 \(\beta\)，定义另一个代数 \({\mathcal A}_\beta\)。

挠 \(H\) 场背景中的弦这一例子，以及相伴的代数 \({\mathcal A}_W\)、\({\mathcal A}_{W'}\)，提示了 \({\mathcal A}_\alpha\) 与 \({\mathcal A}_\beta\) 之间关系的性质：它们是不同的代数，但彼此 Morita 等价，因而具有相同的 \(K\)-理论。

\({\mathcal A}_\alpha\not={\mathcal A}_\beta\) 这一断言如果成立，是相当麻烦的。它是对开弦场论缺乏显式背景独立性的一则尖锐陈述。（此语境中的背景独立性指的是与开弦背景无关；当然，经典开弦场论依赖于开弦在其中传播的那个闭弦背景。）它意味着，这一形式体系依赖于人们在建立理论时使用的是哪个开弦背景。

\(H\) 场中开弦的例子提示了一个补救办法：把 \(D9\)- 与 \(\overline{D9}\)-膜的数目取为无穷。

用 \(m\alpha\) 与 \(m\beta\) 表示 \(\alpha\) 或 \(\beta\) 的 \(m\) 份拷贝（即在 \(\alpha\) 或 \(\beta\) 上补充以 \(U(m)\) 取值的 Chan-Paton 因子）。那么猜想是：对任何有限的 \(m\) 都有 \({\mathcal A}_{m\alpha}\not={\mathcal A}_{m\beta}\)，但当 \(m=\infty\) 时二者相等。其实，关于 \(m=\infty\) 时发生什么的陈述，可以类比（Bouwknegt 与 Mathai 2000）中的陈述而表述得精确得多。令 \({\mathcal K}\) 为可分 Hilbert 空间 \({\mathcal H}\) 中紧算子所成的代数。那么猜想是：对任何 \(\alpha\) 与 \(\beta\)， \begin{equation}\tag{3.5}\label{eq:3.5}
{\mathcal A}_\alpha\otimes {\mathcal K}={\mathcal A}_\beta\otimes {\mathcal K}.
\end{equation}

这一猜想的主要证据是：按照（Bouwknegt 与 Mathai 2000），在有挠 \(H\) 场存在时，相应的陈述（\({\mathcal A}_W\not={\mathcal A}_{W'}\) 但 \({\mathcal A}_W\otimes{\mathcal K}={\mathcal A}_{W'}\otimes {\mathcal K}\)）对零模代数是成立的。此外，如果必须在代数变得同构之前先与 \({\mathcal K}\) 作张量积，这就有助于解释为什么开弦场论中的背景独立性如此难以理解；换言之，它有助于解释为什么用经典解 \(\alpha\) 与代数 \({\mathcal A}_\alpha\) 构造的理论，看起来不同于用经典解 \(\beta\) 与代数 \({\mathcal A}_\beta\) 构造的理论。然而，\SourcePage{13}这一猜想的证据相当有限；也可能只要 \(\alpha\) 与 \(\beta\) 连续连通，就有 \({\mathcal A}_\alpha={\mathcal A}_\beta\)，而理解开弦场论中背景独立性的困难“只不过”是一个技术问题。

如果猜想是对的，人们会猜测：要想在开弦场论中获得更大程度的背景独立性，我们就应当从无穷多个 \(D9\)-膜出发，并依靠快子凝聚把我们带到一个合理的境地。

为使这一点有用，我们需要对 \({\mathcal A}_\alpha\otimes{\mathcal K}\) 有一种描述，它比我们今天对 \({\mathcal A}_\alpha\) 或 \({\mathcal A}_\beta\) 的任何描述都简单得多、也犀利得多。在这样一种假想的新描述中，BRST 算子 \(Q\) 也许会变得更难描述。若果真如此，那我们也只好认了！

\section{数学细节}

这一收尾的小节将用来解释若干我们此前一带而过的数学要点。

设 \({\mathcal A}\) 是一个环——事实上，我们暂时假设 \({\mathcal A}\) 是交换且结合的。若一个 \({\mathcal A}\)-模 \(M\) 是若干个 \({\mathcal A}\) 的直和：\(M={\mathcal A}\oplus {\mathcal A}\oplus \dots\oplus {\mathcal A}\)，则称 \(M\) 为“自由” \({\mathcal A}\)-模。若存在另一个 \({\mathcal A}\)-模 \(M'\)，使得 \(M\oplus M'\) 是自由的，则称 \({\mathcal A}\)-模 \(M\) 为“投射” \({\mathcal A}\)-模。

我们来看当 \({\mathcal A}={\mathcal A}(X)\) 是流形 \(X\) 上连续复值函数的环时，这些定义意味着什么。可以把 \({\mathcal A}\) 看作平凡复线丛 \({\mathcal O}\) 的截面空间。给定一个平凡秩 \(n\) 复向量丛 \(V={\mathcal O}\oplus {\mathcal O}\oplus \dots\oplus{\mathcal O}\)，\(V\) 的截面空间是一个自由模 \(M\)，它是 \(n\) 个 \({\mathcal A}\) 之和。

现在假设给定 \(X\) 上任意一个复向量丛 \(E\)。\(E\) 的截面空间是一个 \({\mathcal A}\)-模 \(M(E)\)（给定函数 \(a\in {\mathcal A}(X)\) 与 \(E\) 的截面 \(m\)，我们直接把 \(a(m)\) 定义为乘积 \(am\)）。（脚注：按我们讨论的非正式风格，我通常对丛和相应的模使用同一个名字 \(E\)。仅在这一段中，我在记号上把 \(E\) 与 \(M(E)\) 区分开来，以便于陈述 Serre-Swan 定理。）\(M(E)\) 是自由模，当且仅当 \(E\) 是平凡向量丛。但 \(M(E)\) 总是投射的。确实，总存在一个“补”向量丛 \(F\)，使得 \(E\oplus F\) 平凡，从而 \(M(E)\oplus M(F)=M(E\oplus F)\) 是自由的。反之（由 Serre 与 Swan 的定理），每个投射 \({\mathcal A}(X)\)-模都是某个 \(X\) 上复向量丛 \(E\) 的 \(M(E)\)。

对任意环 \({\mathcal A}\)，\SourcePage{14}投射模构成一个半群：若 \(E\) 与 \(E'\) 是投射模，则 \(E\oplus E'\) 也是。给定任意半群（其加法运算我们记为 \(\oplus\)），有一种典范的方式构造一个群。这与从正整数半群出发构造整数群的方式相同。群的一个元素是一对 \((E,F)\)，其中 \(E\) 与 \(F\) 是半群的元素（故在我们的情形中，\(E\) 与 \(F\) 是投射模），并服从等价关系 \((E,F)\cong (E\oplus G,F\oplus G)\)，对半群中的任意元素 \(G\) 成立。对的加法为 \((E,F)+(E',F')=(E\oplus E',F\oplus F')\)。等价类构成一个群（零元素或加法单位元是 \(0=(G,G)\)，对任意 \(G\) 成立；\((E,F)\) 的加法逆元是 \((F,E)\)，因为 \((E,F)+(F,E)=(E\oplus F,E\oplus F)=0\)）。若从环 \({\mathcal A}\) 的投射模半群出发，这样构造出的群称为 \(K({\mathcal A})\)。

对于 \({\mathcal A}={\mathcal A}(X)\)，我们在第 1 节中仅用向量丛——换言之，用 \(D9\)-膜与 \(\overline {D9}\)-膜构型——定义了 \(K({\mathcal A})\)。现在我们看到如何以一种具有更广泛有效性的方式来描述这一限制：九维膜对应于投射模，而用九维膜给出的 \(K(X)\) 定义，只是对任意环 \({\mathcal A}\) 用投射模给出 \(K({\mathcal A})\) 定义的一个特例。

现在让我们放弃 \({\mathcal A}\) 交换的假设（但保留结合性）。若 \({\mathcal A}\) 是非交换环，可以典范地给它联系第二个环 \({\mathcal A}^{op}\)（称为相反环），其元素与 \({\mathcal A}\) 的元素一一对应，但乘法顺序相反。于是对每个元素 \(a\in {\mathcal A}\)，有相应的 \(a^{op}\in {\mathcal A}^{op}\)，\({\mathcal A}^{op}\) 的乘法法则为 \(a^{op}b^{op}=(ba)^{op}\)。

对于非交换环，我们必须区分两类模：左模与右模。\({\mathcal A}\) 从左作用于左模 \(M\)，其基本公理是：对 \(a,a'\in {\mathcal A}\) 与 \(m\in M\)，有 \((aa')(m)=a(a'(m))\)。\({\mathcal A}\) 从右作用于右模 \(M\)，其基本公理是 \(m\cdot (aa')=(m\cdot a)\cdot a'\)。\({\mathcal A}\) 的左模与 \({\mathcal A}^{op}\) 的右模相同，反之亦然。

投射模的定义与交换情形相同：左或右 \({\mathcal A}\)-模 \(M\) 是投射的，若存在另一个左或右 \({\mathcal A}\)-模 \(N\)，使得 \(M\oplus N\) 是自由的左或右 \({\mathcal A}\)-模。由投射左 \({\mathcal A}\)-模构成的半群，我们构造出一个群，称为 \(K({\mathcal A})\)。我们也可以由投射右 \({\mathcal A}\)-模构成半群；它与投射左 \({\mathcal A}^{op}\)-模的半群相同，故相应的阿贝尔群是 \(K({\mathcal A}^{op})\)。对于我们关心的环，一般存在一个复（或厄米）共轭运算把 \({\mathcal A}\) 映到 \({\mathcal A}^{op}\)。因此，用左 \({\mathcal A}\)-模或右 \({\mathcal A}\)-模构造的 \(K\)-群\SourcePage{15}在复共轭下彼此等价。

给定两个环 \({\mathcal A}\) 与 \({\mathcal B}\)，一个双模（或 \(({\mathcal A},{\mathcal B})\)-双模）是一个群 \(M\)，它同时是左 \({\mathcal A}\)-模与右 \({\mathcal B}\)-模，且 \({\mathcal A}\) 的左作用与 \({\mathcal B}\) 的右作用彼此交换。对 \(a\in {\mathcal A}\)、\(b\in {\mathcal B}\)、\(m\in M\)，\({\mathcal A}\) 的左作用记为 \(a(m)=am\)，\({\mathcal B}\) 的右作用记为 \(b(m)=mb\)；由于我们假设 \({\mathcal A}\) 与 \({\mathcal B}\) 的作用交换，有 \((am)b=a(mb)\)；我们将这些表达式简写为 \(amb\)。

现在我们可以定义环之间一个非常重要的关系，称为 Morita 等价\TN{Morita equivalence；旧译名不统一，本文一律作 Morita 等价。}。它最早是在 [20] 中与 \(D\)-膜联系起来使用的；进一步的进展见（Schwarz 1998；Schwarz 与 Pioline 1999）。设 \(M\) 是一个 \(({\mathcal A},{\mathcal B})\)-双模。\(M\) 的选取使我们能定义一个从左 \({\mathcal B}\)-模到左 \({\mathcal A}\)-模的映射：给定左 \({\mathcal B}\)-模 \(W\)，把它映为左 \({\mathcal A}\)-模 \(W'=M\otimes_{{\mathcal B}}W\)。类似地，若给定一个 \(({\mathcal B},{\mathcal A})\)-双模 \(N\)，我们可以通过 \(W^{\prime\prime}=N\otimes_{\mathcal A}W'\) 把左 \({\mathcal A}\)-模 \(W'\) 映为左 \({\mathcal B}\)-模 \(W^{\prime\prime}\)。如果这两个运算互为逆运算，换言之，若对所有 \(W\) 都有 \(W^{\prime\prime}=W\)，我们就说 \({\mathcal A}\) 与 \({\mathcal B}\) 是 Morita 等价的。此时，左 \({\mathcal A}\)-模的半群同构于左 \({\mathcal B}\)-模的半群，且 \(K({\mathcal A})=K({\mathcal B})\)。

这一关系在开弦场论中的含义已在（Seiberg 与 Witten 1999）第 6.4 节中解释。对任意开弦边界条件 \(\alpha\)，可定义一个由 \(\alpha\)-\(\alpha\) 开弦构成的环 \({\mathcal A}_{\alpha}\)。对任意其他边界条件 \(\gamma\)，\(\alpha\)-\(\gamma\) 开弦构成一个左 \({\mathcal A}_\alpha\)-模 \(M_{\alpha\gamma}\)，而 \(\gamma\)-\(\alpha\) 开弦构成一个右 \({\mathcal A}_\alpha\)-模 \(M_{\gamma\alpha}\)。想必在适当的意义上，它们构成左、右 \({\mathcal A}_\alpha\)-模的一组基。现在考虑任意其他开弦边界条件 \(\beta\) 及相应的代数 \({\mathcal A}_\beta\)。用 \(\beta\)-\(\gamma\) 与 \(\gamma\)-\(\beta\) 开弦，以同样的方式构造左、右 \({\mathcal A}_\beta\)-模 \(M_{\beta\gamma}\) 与 \(M_{\gamma\beta}\)。显然，左 \({\mathcal A}_\alpha\)-模与左 \({\mathcal A}_\beta\)-模之间存在自然的对应，即

\begin{equation}\tag{4.1}\label{eq:4.1}
M_{\alpha\gamma}\leftrightarrow M_{\beta\gamma}.
\end{equation}

因此对所有 \(\alpha\) 与 \(\beta\) 都有 \(K({\mathcal A}_\alpha)=K({\mathcal A}_\beta)\)。我们此前不借助 Morita 等价直接确立了后一事实，但实际上，\({\mathcal A}_\alpha\) 与 \({\mathcal A}_\beta\) 之间的关系似乎正是它们是 Morita 等价的。确实，由 \(\alpha\)-\(\beta\) 与 \(\beta\)-\(\alpha\) 开弦可以构造出自然的 \(({\mathcal A}_\alpha,{\mathcal A}_\beta)\) 与 \(({\mathcal A}_\beta,{\mathcal A}_\alpha)\) 双模 \(M_{\alpha\beta}\) 与 \(M_{\beta\alpha}\)；这些双模极有可能具备确立 \({\mathcal A}_\alpha\) 与 \({\mathcal A}_\beta\) 之间 Morita 等价所需的性质。

现在我们考虑这一点的更初等版本，目的是为第 3.2 节的讨论补充一些细节。设 \(M_n\) 是 \(n\times n\) 复矩阵的代数。\SourcePage{16}最显然的左 \(M_n\)-模是 \(n\) 分量列向量的空间 \(W\)。这是一个投射模，因为作为左模，\(M_n\) 由 \(n\) 个 \(W\) 的拷贝组成（\(M_n\) 由 \(n\times n\) 矩阵组成；在 \(M_n\) 对自身的左作用中，\(M_n\) 分别作用于构成 \(n\times n\) 矩阵的 \(n\) 列中的每一列，故这些列构成 \(n\) 个 \(W\) 的拷贝）。所以若 \(W'\) 是 \(n-1\) 个 \(W\) 之和，则 \(W\oplus W'\) 是自由的，因而 \(W\) 是投射的。可以证明，每个投射左 \(M_n\)-模都是若干个 \(W\) 之和。于是投射左 \(M_n\)-模的半群同构于非负整数半群，且 \(K(M_n)=\mathbb{Z}\)。

既然 \(K(M_n)\) 因此与 \(n\) 无关，人们或许会问，不同 \(n\) 的各个 \(M_n\) 是否 Morita 等价。确实如此。对任意正整数 \(n\) 与 \(k\)，最显然的 \((M_n,M_k)\)-双模是 \(n\times k\) 矩阵的空间 \(C_{n,k}\)，其中 \(M_n\) 从左作用、\(M_k\) 从右作用。同样，最显然的 \((M_k,M_n)\)-双模是 \(C_{k,n}\)。可以证明，这些双模确立了 \(M_n\) 与 \(M_k\) 之间的 Morita 等价。

现在，我们重新考虑第 3.2 节讨论过的挠 \(H\)-场。对任意扭曲丛 \(W\)，存在一个由 \(W\)-\(W\) 开弦的基态生成的代数 \({\mathcal A}_W(X)\)。\({\mathcal A}_W(X)\) 的元素是 \(W\otimes \overline W\) 的截面。给定任意两个不同的扭曲丛 \(W\) 与 \(W'\)，我们可以把上面关于一般开弦代数 \({\mathcal A}_\alpha\) 与 \({\mathcal A}_\beta\) 的更抽象讨论具体应用到这一情形，从而试着在 \({\mathcal A}_W(X)\) 与 \({\mathcal A}_{W'}(X)\) 之间建立 Morita 等价。这意味着我们应当通过考察 \(W\)-\(W'\) 与 \(W'\)-\(W\) 开弦的基态，来找到 \(({\mathcal A}_W(X),{\mathcal A}_{W'}(X))\) 与 \(({\mathcal A}_{W'}(X),{\mathcal A}_W(X))\) 双模。\(W\)-\(W'\) 开弦的基态是 \(W\otimes \overline{W'}\) 的截面；这些截面的空间是一个 \(({\mathcal A}_W(X),{\mathcal A}_{W'}(X))\) 双模 \(M\)。同样，\(W'\)-\(W\) 开弦基态的空间是一个由 \(W'\otimes \overline W\) 的截面构成的 \(({\mathcal A}_{W'}(X),{\mathcal A}_W(X))\) 双模 \(N\)。双模 \(M\) 与 \(N\) 确实在 \({\mathcal A}_W(X)\) 与 \({\mathcal A}_{W'}(X)\) 之间确立了 Morita 等价。确实，在证明这一点时可以在 \(X\) 上局部地进行。局部地，\(W\) 与 \(W'\) 是平凡的，并且（若 \(n\) 与 \(k\) 分别是 \(W\) 与 \(W'\) 的秩）如第 3.2 节所述，局部地有 \({\mathcal A}_W(X)={\mathcal A}(X)\otimes M_n\)、\({\mathcal A}_{W'}(X)={\mathcal A}(X)\otimes M_k\)。给出 \(M_n\) 与 \(M_k\) 的 Morita 等价的同一计算表明，\({\mathcal A}_W(X)\) 与 \({\mathcal A}_{W'}(X)\) 是 Morita 等价的。

投射右 \({\mathcal A}_W(X)\) 模是 \(E\otimes \overline W\) 的截面空间，其中 \(E\) 是任意扭曲丛。（若 \(s^\alpha{}_j\) 是这样一个截面——其中 \(\alpha\) 是 \(E\) 指标、\(j\) 是 \(\overline W\) 指标——而 \(w^k{}_l\) 是 \(W\otimes \overline W\) 的截面，则 \(s\) 与 \(w\) 的乘积定义为 \((sw)^\alpha{}_j=\sum_ks^\alpha{}_kw^k{}_j\)。）\SourcePage{17}所以投射右 \({\mathcal A}_W(X)\) 模与扭曲丛 \(E\) 自然地一一对应，而投射右 \({\mathcal A}_W(X)\) 模的 \(K\)-群就是受通常等价关系约束的对 \((E,F)\) 构成的群。对任意两个具有相同 \(H\) 的扭曲丛 \(W\) 与 \(W'\)，\({\mathcal A}_W\) 与 \({\mathcal A}_{W'}\) 的 Morita 等价只是以一种更讲究的方式，描述由对应

\begin{equation}\tag{4.2}\label{eq:4.2}
E\otimes \overline W \leftrightarrow E\otimes \overline{W'}
\end{equation}

给出的 \({\mathcal A}_W\)-模与 \({\mathcal A}_{W'}\)-模之间的自然映射。右 \({\mathcal A}_W\)-模与右 \({\mathcal A}_{W'}\)-模之间的这个一一映射（或左模的相应映射）导致 \(K({\mathcal A}_W)=K({\mathcal A}_{W'})\) 的结论。

本工作部分得到 NSF 基金 PHY-9513835 和 Caltech Discovery Fund 的资助。感谢 K. Hori 与 G. B. Segal 的评论与解释。

\SourcePage{18}
\section*{参考文献}
\noindent{\footnotesize 译注：参考文献保留原文英文条目，编号与原文一致；未逐条中译，以免二次转写失真。}\par\smallskip
{\small
\begin{enumerate}[label={[\arabic*]},leftmargin=2.4em,itemsep=2pt,topsep=4pt]
\item A. Sen, ``Tachyon Condensation On The Brane-Antibrane System,'' JHEP 9808:012 (1998), hep-th/9805170.
\item E. Witten, ``D-Branes And K-Theory,'' JHEP 9812:019 (1998), hep-th/9810188.
\item E. Diaconescu, G. Moore, and E. Witten, ``E8 Gauge Theory, And A Derivation Of K-Theory From M-Theory,'' hep-th/0005091.
\item D. Freed and E. Witten, ``Anomalies In String Theory With D-Branes,'' hep-th/9907189.
\item A. Connes, M. Douglas, and A. Schwarz, ``Noncommutative Geometry And Matrix Theory: Compactification On Tori,'' JHEP 9802:003 (1998), hep-th/9711162.
\item D. S. Freed and M. J. Hopkins, ``On Ramond-Ramond Fields And K-Theory,'' JHEP 0005:044 (2000), hep-th/0002027.
\item E. Witten, ``Duality Relations Among Topological Effects In String Theory,'' JHEP 0005:031 (2000), hep-th/9912086.
\item G. Moore and E. Witten, ``Self-Duality, Ramond-Ramond Fields, And K-Theory,'' JHEP 0005:032 (2000), hep-th/9912279.
\item M. Bershadsky, C. Vafa, and V. Sadov, ``D-Branes And Topological Field Theories,'' Nucl. Phys. B463 (1996) 420.
\item M. B. Green, J. A. Harvey, and G. Moore, ``I-Brane Inflow And Anomalous Couplings On D-Branes,'' Class. Quant. Grav. 14 (1997) 47, hep-th/9605033.
\item E. Cheung and Z. Yin, ``Anomalies, Branes, and Currents,'' Nucl. Phys. B517 (1998) 69, hep-th/9710206.
\item G. Moore and R. Minasian, ``K-Theory And Ramond-Ramond Charge,'' JHEP 9711:002 (1997), hep-th/9710230.
\item E. Witten, ``On Flux Quantization In M-Theory And The Effective Action,'' J. Geom. Phys. 22 (1997) 1, hep-th/9609122.
\item M. Fabinger and P. Horava, ``Casimir Effect Between World-Branes In Heterotic M-Theory,'' Nucl. Phys. B580 (2000) 243, hep-th/0002073.
\item K. Hori, ``D-Branes, T-Duality, and Index Theory,'' hep-th/9902102.
\item P. Horava, ``Type IIA D-Branes, K Theory, and Matrix Theory,'' ATMP 2 (1999) 1373, hep-th/9812135.
\item A. Kapustin, ``D-Branes In A Topologically Nontrivial B-Field,'' hep-th/9909089.
\item P. Bouwknegt and V. Mathai, ``D-Branes, B Fields, And Twisted K Theory,'' JHEP 0003:007 (2000), hep-th/0002023.
\item M. F. Atiyah and G. B. Segal, unpublished.
\item A. Schwarz, ``Morita Equivalence And Duality,'' Nucl. Phys. B534 (1998) 720, hep-th/9805034; A. Schwarz and B. Pioline, ``Morita Equivalence And T Duality (Or B Versus Theta),'' JHEP 9908:021 (1999), hep-th/9908019.
\item N. Seiberg and E. Witten, ``String Theory And Noncommutative Geometry,'' JHEP 9909:032 (1999), hep-th/9908142.
\end{enumerate}
}

\end{document}
